% ============================================================
\chapter{Fourier Transform and PDEs}
\label{ch:fourier_transform}
% ============================================================

In 1807, Joseph Fourier claimed that every function can be written as a sum of sinusoids. Lagrange was sceptical, Poisson doubtful, but the idea was too powerful to be ignored. The \emph{Fourier transform}, the continuous generalization of Fourier series, converts a function of the spatial variable into a function of frequency. For PDEs with constant coefficients, this transformation is miraculous: derivatives become multiplications, partial differential equations become algebraic equations. The heat equation, the wave equation, the Schr\"odinger equation --- all are solved elegantly by Fourier analysis on $\R^n$. This chapter develops the Fourier transform on $L^1$ and $L^2$, establishes the inversion formula and Plancherel's theorem, then systematically applies them to classical PDEs.

\section{Fourier transform on $L^1(\R^n)$}

\begin{definition}[Fourier transform]
\index{Fourier transform}
Let $f \in L^1(\R^n)$. The \emph{Fourier transform} of $f$ is the function
$\hat{f} : \R^n \to \C$ defined by:
\[
    \hat{f}(\xi) = \mathcal{F}[f](\xi) = \int_{\R^n} f(x)\, e^{-2\pi i \inner{x}{\xi}} \, dx.
\]
\end{definition}

\begin{remark}
There are several conventions for the Fourier transform. We use the symmetric
convention with $2\pi$ in the exponential, which eliminates factors of
$(2\pi)^{n/2}$ from the inversion formula.
\end{remark}

\begin{proposition}[Basic properties]
\label{prop:fourier_properties}
Let $f, g \in L^1(\R^n)$:
\begin{enumerate}
    \item \textbf{Linearity}: $\widehat{\alpha f + \beta g} = \alpha \hat{f} + \beta \hat{g}$.
    \item \textbf{Translation}: $\widehat{f(\cdot - a)}(\xi) = e^{-2\pi i \inner{a}{\xi}} \hat{f}(\xi)$.
    \item \textbf{Modulation}: $\widehat{e^{2\pi i \inner{a}{\cdot}} f}(\xi) = \hat{f}(\xi - a)$.
    \item \textbf{Dilation}: $\widehat{f(\lambda \cdot)}(\xi) = \frac{1}{\abs{\lambda}^n} \hat{f}(\xi/\lambda)$ for $\lambda \neq 0$.
    \item \textbf{Convolution}: $\widehat{f * g} = \hat{f} \cdot \hat{g}$.
\end{enumerate}
\end{proposition}

\begin{theorem}[Riemann-Lebesgue lemma]
\index{Riemann-Lebesgue lemma}
If $f \in L^1(\R^n)$, then $\hat{f} \in C_0(\R^n)$ (continuous and vanishing at
infinity): $\hat{f}(\xi) \to 0$ as $\abs{\xi} \to \infty$.
\end{theorem}

\begin{proof}
For $f$ the indicator of a rectangle, the result is a direct computation. By
density of step functions in $L^1$ and continuity of $\mathcal{F}$ (since
$\abs{\hat{f}(\xi)} \leq \norm{f}_{L^1}$), the result extends to all $f \in L^1$.
\end{proof}

\begin{theorem}[Transform of derivatives]
\index{Fourier transform!derivatives}
If $f$ and $\partial^\alpha f$ are in $L^1(\R^n)$, then:
\[
    \widehat{\partial^\alpha f}(\xi) = (2\pi i \xi)^\alpha \hat{f}(\xi),
\]
where $\alpha$ is a multi-index and $(2\pi i \xi)^\alpha = \prod_{j=1}^{n}(2\pi i \xi_j)^{\alpha_j}$.
\end{theorem}

\begin{intuition}
The Fourier transform converts derivatives into multiplication by polynomials.
This is the fundamental reason for its effectiveness in solving linear PDEs with
constant coefficients: a PDE becomes an algebraic equation in Fourier space.
\end{intuition}

\section{Fourier transform on $L^2(\R^n)$}

\begin{theorem}[Plancherel]
\label{thm:plancherel}
\index{Plancherel theorem}
The Fourier transform extends to an isometric isomorphism of $L^2(\R^n)$:
\[
    \norm{\hat{f}}_{L^2} = \norm{f}_{L^2} \quad \forall\, f \in L^2(\R^n).
\]
More generally, $\inner{\hat{f}}{\hat{g}}_{L^2} = \inner{f}{g}_{L^2}$
(Parseval's identity).
\end{theorem}

\begin{proof}[Proof sketch]
One first proves the result for $f \in L^1 \cap L^2$ (a dense subset of $L^2$)
by computing $\int \abs{\hat{f}}^2$ via the convolution
$f * \overline{f(-\cdot)}$. The result then extends by density and continuity.
\end{proof}

\begin{theorem}[Inversion formula]
\index{Fourier inversion formula}
If $f, \hat{f} \in L^1(\R^n)$, then:
\[
    f(x) = \int_{\R^n} \hat{f}(\xi)\, e^{2\pi i \inner{x}{\xi}} \, d\xi
    \quad \text{a.e.}
\]
\end{theorem}

\section{Tempered distributions}

\begin{definition}[Schwartz space]
\index{Schwartz space}
The \emph{Schwartz space} $\mathcal{S}(\R^n)$ is the space of functions
$\varphi \in C^\infty(\R^n)$ whose derivatives all decay faster than any
polynomial:
\[
    \sup_{x \in \R^n} \abs{x^\alpha \partial^\beta \varphi(x)} < \infty
    \quad \forall\, \alpha, \beta \in \N^n.
\]
\end{definition}

\begin{definition}[Tempered distribution]
\index{tempered distribution}
The space $\mathcal{S}'(\R^n)$ of \emph{tempered distributions} is the
topological dual of $\mathcal{S}(\R^n)$. The Fourier transform extends to
$\mathcal{S}'$ by duality:
\[
    \langle \hat{T}, \varphi \rangle = \langle T, \hat{\varphi} \rangle
    \quad \forall\, T \in \mathcal{S}', \; \varphi \in \mathcal{S}.
\]
\end{definition}

\begin{example}[Fourier transform of the Dirac mass]
$\hat{\delta} = 1$, since $\langle \hat{\delta}, \varphi \rangle = \langle \delta, \hat{\varphi} \rangle
= \hat{\varphi}(0) = \int \varphi(\xi)\, d\xi = \langle 1, \varphi \rangle$.
\end{example}

\begin{example}[Fourier transform of the Gaussian]
\index{Gaussian}
For $f(x) = e^{-\pi \abs{x}^2}$, we have $\hat{f}(\xi) = e^{-\pi \abs{\xi}^2}$.
The Gaussian is a fixed point of the Fourier transform.
\end{example}

\section{Solving the heat equation}

Consider the Cauchy problem for the heat equation:
\begin{equation}\label{eq:heat_cauchy}
    u_t = \kappa\, \Delta u, \quad x \in \R^n,\; t > 0, \qquad
    u(x, 0) = u_0(x).
\end{equation}

\begin{theorem}[Solution via Fourier transform]
\label{thm:heat_fourier}
The solution of \eqref{eq:heat_cauchy} is given by:
\[
    u(x, t) = (G_t * u_0)(x) = \int_{\R^n} G_t(x - y)\, u_0(y)\, dy,
\]
where $G_t$ is the \emph{heat kernel}:
\[
    G_t(x) = \frac{1}{(4\pi \kappa t)^{n/2}} \exp\!\left(-\frac{\abs{x}^2}{4\kappa t}\right).
\]
\end{theorem}

\begin{proof}
Applying the Fourier transform in $x$:
\[
    \hat{u}_t(\xi, t) = -4\pi^2 \kappa \abs{\xi}^2 \hat{u}(\xi, t).
\]
This is an ODE in $t$, with solution
$\hat{u}(\xi, t) = \hat{u}_0(\xi)\, e^{-4\pi^2 \kappa \abs{\xi}^2 t}$.
Recognizing $e^{-4\pi^2 \kappa \abs{\xi}^2 t} = \hat{G}_t(\xi)$, we obtain
$u = G_t * u_0$ by the convolution property.
\end{proof}

\begin{warning}[title=\textbf{Instantaneous propagation}]
The heat kernel satisfies $G_t(x) > 0$ for all $x$ and all $t > 0$. Thus, even
if $u_0$ has compact support, $u(x, t) > 0$ for all $x$ as soon as $t > 0$:
heat propagates at infinite speed.
\end{warning}

\section{Solving the wave equation}

\begin{theorem}[Solution of the wave equation in dimension $n$]
The Cauchy problem:
\[
    u_{tt} = c^2 \Delta u, \quad u(x,0) = u_0(x), \quad u_t(x,0) = v_0(x),
\]
has the solution (in Fourier space):
\[
    \hat{u}(\xi, t) = \hat{u}_0(\xi) \cos(2\pi c \abs{\xi} t)
    + \hat{v}_0(\xi) \frac{\sin(2\pi c \abs{\xi} t)}{2\pi c \abs{\xi}}.
\]
\end{theorem}

\begin{proof}
The Fourier transform in $x$ gives
$\hat{u}_{tt} = -(2\pi c)^2 \abs{\xi}^2 \hat{u}$. For each fixed $\xi$, this
is a harmonic oscillator equation with frequency $\omega = 2\pi c \abs{\xi}$.
The general solution is $\hat{u}(\xi, t) = A(\xi)\cos(\omega t) + B(\xi)\sin(\omega t)$.
Initial conditions give $A = \hat{u}_0$ and $B = \hat{v}_0 / \omega$.
\end{proof}

\begin{example}[d'Alembert's formula ($n = 1$)]
In dimension $1$, inverting gives:
\[
    u(x,t) = \frac{u_0(x+ct) + u_0(x-ct)}{2} + \frac{1}{2c}\int_{x-ct}^{x+ct} v_0(s)\, ds.
\]
\end{example}

\section{Fundamental solutions}

\begin{definition}[Fundamental solution]
\index{fundamental solution}
The \emph{fundamental solution} (or Green's function) of a linear operator $L$
is the distribution $E$ such that $LE = \delta$.
\end{definition}

\begin{example}[Laplacian in dimension $n \geq 3$]
The fundamental solution of $-\Delta$ in $\R^n$ ($n \geq 3$) is:
\[
    E(x) = \frac{1}{n(n-2)\omega_n} \frac{1}{\abs{x}^{n-2}},
\]
where $\omega_n = \frac{2\pi^{n/2}}{\Gamma(n/2)}$ is the area of $S^{n-1}$.
\end{example}

\begin{keyformulas}[title=\textbf{Fundamental solutions of classical PDEs}]
\begin{align*}
    \text{Heat:} \quad & G_t(x) = \frac{1}{(4\pi \kappa t)^{n/2}}
        e^{-\abs{x}^2/(4\kappa t)}, \quad t > 0 \\[6pt]
    \text{Laplacian ($n=3$):} \quad & E(x) = \frac{1}{4\pi \abs{x}} \\[6pt]
    \text{Laplacian ($n=2$):} \quad & E(x) = -\frac{1}{2\pi}\ln\abs{x} \\[6pt]
    \text{Helmholtz ($n=3$):} \quad & E(x) = \frac{e^{ik\abs{x}}}{4\pi \abs{x}}
\end{align*}
\end{keyformulas}

\section{Exercises}

\begin{exercise}
Compute the Fourier transform of $f(x) = e^{-a\abs{x}}$ for $a > 0$ and $x \in \R$.
\end{exercise}

\begin{exercise}
Prove the convolution property: if $f, g \in L^1(\R^n)$, then
$\widehat{f * g} = \hat{f} \cdot \hat{g}$.
\end{exercise}

\begin{exercise}
Verify that the heat kernel $G_t$ satisfies:
(a) $\int_{\R^n} G_t(x)\, dx = 1$ for all $t > 0$;
(b) $(\partial_t - \kappa \Delta)G_t = 0$ for $t > 0$;
(c) $G_t \to \delta$ in the sense of distributions as $t \to 0^+$.
\end{exercise}

\begin{exercise}
Use the Fourier transform to solve the free Schr\"odinger equation:
$i\hbar \psi_t = -\frac{\hbar^2}{2m}\Delta \psi$, $\psi(x,0) = \psi_0(x)$.
\end{exercise}

\begin{exercise}
Let $f \in \mathcal{S}(\R^n)$. Show that if $\hat{f}$ has compact support, then
$f$ extends to an entire function on $\C^n$ (Paley-Wiener theorem).
\end{exercise}

\begin{exercise}
Solve the heat equation on $\R$ with initial data $u_0(x) = \ind_{[0,1]}(x)$
(indicator function). Express the solution using the error function $\operatorname{erf}$.
\end{exercise}

\begin{exercise}
Show that the fundamental solution of the Laplacian in $\R^n$ ($n \geq 3$)
satisfies $-\Delta E = \delta$ in the sense of distributions, by computing
$\Delta E$ for $x \neq 0$ and using the divergence theorem.
\end{exercise}
